\documentclass{article}
\usepackage{graphicx}
\usepackage[margin=1in]{geometry}
\usepackage{tikz} % for diagram
\usetikzlibrary{positioning}
\usepackage{hyperref}
\hypersetup{
    colorlinks=true,
    linkcolor= dark blue,
    filecolor=magenta,      
    urlcolor=cyan,
    pdftitle={Overleaf Example},
    pdfpagemode=FullScreen,
    }


\title{\bf{Project:\\ Documentation}}
% \author{Monem Moheq}
\date{June 2026}

\begin{document}

\maketitle

\section*{Data Collection}

    \subsection*{Calgary Data Attempt}

         \begin{flushleft}
            \textbf{6/23/2026}
        \end{flushleft}
    
        We use the \href{https://data.calgary.ca/Business-and-Economic-Activity/Building-Permits/c2es-76ed/about_data}{Calgary Data Base}. Through the online querying system, we have filtered the data as follows.

        \begin{itemize}
            \item Permit type: Commercial/multi family
            \item Workclassgroup: new
            \item Applied date: 2017 Jun 22 to 2023 Jun 22
            \item PermitClassMapped: Residential
            \item StatusCurrent: Complete
            \item PermitClassGroup: Unspecified
        \end{itemize}
        
        Issue: Some projects show up as two pieces. Ie, the complex shows up as two separate rows even though it may be considered one project. We are considering ignoring this, treat them as separate projects, or we can aggregate and combine them. 
        
        We use the assessed value at the year the building permit is obtained as the initial price. We use the assessed value at the year of completion as the final price. The assessed values must be obtained from this from the \href{https://data.calgary.ca/Government/Total-Property-Assessed-Value/dmd8-bmxh}{assessed value data base}. We have downloaded the data for the year 2017 to 2023, it is on Finley's computer, it is a larger file about 2.5GB.

        \begin{flushleft}
            \textbf{6/27/2026}
        \end{flushleft}

        To expound upon the existing data, since aggregation of building permits belonging to the same projects produced a smaller data set than anticipated, we have expanded to data spanning 2013-2023. Both the 2013 and 2017 scenarios have been pinned on the discord. The datasets, as of right now, are based off of townhouse projects. We can look at apartment projects, but those data points seem a little more messy as some non-apartment permits are sprinkled throughout but only as differentiated as non-apartment in their text descriptions, not by any organized category. 

        We have begun considering using multiple data sets from different cities, such as Vancouver, Burnaby, etc., and having the city be another factor in the prediction. So long as we can gather good data from these places, this should be a feature to implement and will increase the versatility of the model. 

        As of right now, working on getting a basic classifier running (or even the random forest as well, although we will have to retrain it later) would be good as we can work on increasing the data set concurrently, as we can test the basic functionality of the model using these smaller datasets.
\\\\
        \begin{flushleft}
            \textbf{6/29/2026}
        \end{flushleft}

        We wish to match the assessed value of properties data from one dataset, to the town house properties previously narrowed down in the other dataset. For this we use a python script and the pandas library to clean and combine the data. This has given rise to a new problem. After combining the data, 68 properties are left unmatched. This could be the result of poor execution of the combining of data, or simply that the assessed value data set does not have a value for some properties. Further investigation is required. The cleaning and combination of the data ran on some assumptions that also need to be investigated. Some addresses in the assessment dataset appear multiple times. In cleaning the data, duplicate addresses were removed, keeping the first one that appeared, in the default order of the data. This may not be ideal. Another assumption was that the addresses are standardized between the data sets, if this is not the case, it could explain the unmatched assessment values in the combined data. The combined data is uploaded to the discord, as well as the python script. 

        \begin{flushleft}
            \textbf{7/05/2026}
        \end{flushleft}

        Centroid coordinate has been added to each project as well as distance to nearest transit stop and also distance to downtown. consolidating the income data and the project data is difficult. Some community names from the project data match those in the income data, but not all. One must explore how we can connect the project location to the average income of that neighbourhood.
    

        
\end{document}
